\documentclass[12pt]{article}
\usepackage{amsmath}
\usepackage{amssymb}
\usepackage[T1]{fontenc}
\usepackage{xcolor}
\usepackage{listings}
\usepackage{geometry}
\usepackage{graphicx}
\geometry{margin=1in}

\title{Larkos Training Loop and Inference Runner}
\author{Witold Warcho\l{}}
\date{\today}

\lstset{
  basicstyle=\ttfamily\small,
  keywordstyle=\color{blue},
  commentstyle=\color{gray},
  stringstyle=\color{red},
  breaklines=true,
  captionpos=b
}

\begin{document}

\maketitle

\section{Overview}

The Larkos system consists of two major components:

\begin{enumerate}
  \item \textbf{TrainingLoop}: Full differentiable training with gradient
    computation, emotional state updates, and cognitive fusion
  \item \textbf{LarkosRunner}: Inference-only forward pass that loads
    trained weights and runs without gradient accumulation
\end{enumerate}

Both components share the same core modules: the LarkosModel,
\_FusionTransformerHead, embedding networks, and text codec. The key
difference is that TrainingLoop computes gradients and updates weights,
while LarkosRunner only performs forward inference using frozen weights.

\section{Core Mathematical Components}

\subsection{Fourier Encoding}

The input $x \in \mathbb{R}^D$ is encoded using positional-style Fourier
features:

\[
  \phi(x) = \left[ \sin(2^0 \pi x), \sin(2^1 \pi x), \ldots,
    \sin(2^{E-1} \pi x), \cos(2^0 \pi x), \cos(2^1 \pi x), \ldots,
    \cos(2^{E-1} \pi x) \right]
\]

where $E$ is the number of encodings ($\texttt{FOURIER\_ENCODINGS}$). The
output has dimension $2 \cdot D \cdot E$.

\textbf{Rationale}: Fourier encodings help the model learn periodic and
frequency-dependent patterns. They provide stable gradients compared to
raw inputs and help with generalization.

\subsection{Online Normalization}

Two normalization strategies track running statistics:

\subsubsection{OnlineMinMax}

Per-dimension running min/max via exponential moving average (EMA):

\[
  m_t^{(i)} = m_{t-1}^{(i)} + \alpha (x_t^{(i)} - m_{t-1}^{(i)})
\]

where $\alpha = 0.02$ (momentum). The normalized output is:

\[
  \tilde{x}^{(i)} = \frac{x^{(i)} - m}{M - m + \epsilon} \cdot 2 - 1
\]

clamped to $[-3, 3]$.

\textbf{Rationale}: EMA-based min/max prevents early outliers from
permanently warping the normalization range. The system gradually forgets
extreme values and re-centers around the actual distribution.

\subsubsection{OnlineMeanStd}

Per-dimension running mean and variance via EMA:

\[
  \mu_t^{(i)} = \mu_{t-1}^{(i)} + \alpha (x_t^{(i)} - \mu_{t-1}^{(i)})
\]

\[
  \sigma_t^{(i)} = (1 - \alpha) \left( \sigma_{t-1}^{(i)} +
    \alpha (x_t^{(i)} - \mu_{t-1}^{(i)})^2 \right)
\]

Output normalization:

\[
  \hat{x}^{(i)} = \frac{x^{(i)} - \mu}{\sqrt{\sigma} + \epsilon_f}
\]

where $\epsilon_f = 0.1$ (std floor).

\textbf{Rationale}: The std floor prevents near-constant dimensions from
dividing by zero and blowing up. It keeps low-variance dimensions genuinely
quiet rather than amplified by numerical instability.

\subsection{Temporal Context Window}

The Fourier-encoded input is padded into a temporal window of size
$\texttt{TEMPORAL\_WINDOW}$:

\[
  x_{\text{temporal}} = \text{concat}(x_{t-w+1}, x_{t-w+2}, \ldots, x_t)
\]

where $w = \texttt{TEMPORAL\_WINDOW}$. Earlier positions are zero-padded.

\textbf{Rationale}: Temporal history provides context for the model to
detect trends and patterns over time, improving predictions.

\section{Training Loop Architecture}

\subsection{Forward Pass}

The forward pass (\texttt{\_forward}) builds a fully differentiable
computation graph:

\[
  \text{model\_pred} = \text{LarkosModel}(x_{\text{temporal}},
    \text{embed\_ctx})
\]

The LarkosModel outputs a vector that is split into an LLM embedding
($\text{llm\_embed}$) and auxiliary dimensions.

\subsubsection{Cross-Attention Path}

The LLM embedding is gated and fused with text context:

\[
  \text{embed\_gate} = \text{EmbedWeightNet}(x_{\text{norm}})
\]

\[
  \text{llm\_embed\_g} = \text{llm\_embed} \odot \text{embed\_gate}
\]

\[
  \text{llm\_embed\_ca} = \text{InputCrossAttention}(x_{\text{norm}},
    \text{llm\_embed\_g})
\]

where $\odot$ denotes element-wise multiplication.

\textbf{Why Cross-Attention}: Instead of concatenating the input and
embedding, cross-attention treats the $D$-dimensional input as $D$ separate
1-D query tokens and the $E$-dimensional embedding as key/value. This is
far more expressive than simple concatenation and allows learned
interaction.

\subsubsection{Cognitive Fusion}

The cross-attention embedding is detached (breaking the gradient graph at
the C boundary) and passed to cognitive fusion:

\[
  \text{fused\_cog\_raw} = \text{cognitive\_fuse}(\text{llm\_embed\_ca},
    \text{neurons}, \text{mem\_state}, \text{default\_weights}, \ldots)
\]

\textbf{Graph Breaking}: The detach at the C boundary is intentional. The
C-side fusion depends on backend neurons and memory state which are
non-differentiable objects. The Python side cannot backprop into them. An
auxiliary loss on the cross-attention embedding (\texttt{aux\_proj})
provides the gradient signal for the modules before the C boundary.

\subsubsection{Frozen Input Windows}

Training targets are refreshed every $\texttt{TARGET\_FREEZE\_INTERVAL}$
epochs. Between refreshes, both the cached target and cached fused
cognitive vector are held constant:

\[
  \text{frozen\_input} = \text{(cache exists)} \land
    \text{(epoch} - \text{target\_epoch)} < \text{INTERVAL}
\]

\textbf{Rationale}: If the target is fixed but the transformer input
drifts, the model chases a moving input toward a fixed point, producing
unstable loss. Pinning both together and refreshing them synchronously
prevents this pathology. The frozen window allows the model to converge
on the pair without the input changing underneath it.

\subsection{Fusion Transformer Head}

The \_FusionTransformerHead treats the fused cognitive vector as three
contiguous bands that form a 3-token sequence:

\[
  \text{Band 0} = \text{fused}[0:\text{BAND\_Q}] \quad (\text{LLM query
    stream})
\]

\[
  \text{Band 1} = \text{fused}[\text{BAND\_Q}:\text{BAND\_Q+BAND\_N}]
    \quad (\text{neuron stream})
\]

\[
  \text{Band 2} = \text{fused}[\text{BAND\_Q+BAND\_N}:]
    \quad (\text{memory stream})
\]

Each band is projected into a shared $d_{\text{model}}$ space:

\[
  \text{tok}_i = \text{proj}_i(\text{Band}_i) + \text{stream\_embed}_{i}
\]

A transformer encoder runs self-attention over the 3 tokens:

\[
  \text{enc} = \text{TransformerEncoder}([\text{tok}_{0}, \text{tok}_{1},
    \text{tok}_{2}])
\]

Output is mean-pooled and projected to neuron dimensions:

\[
  \text{fused} = \text{LinearOut}\left(
    \frac{1}{3} \sum_{i} \text{enc}_{i} \right)
\]

\textbf{Why Three Bands}: The bands carry different semantic information
(LLM query, neuron state, memory). Using separate projections and stream
embeddings lets each stream learn its own mapping and lets attention
determine their balance.

\textbf{Mean Pooling}: Pooling instead of taking a single token keeps all
three streams in the output. Attention learns to weight each stream
appropriately rather than collapsing to one.

\subsection{Loss Computation}

The loss is a weighted blend of four components:

\[
  L_{\text{total}} = \lambda_1 L_{\text{outer}} + \lambda_2 L_{\text{base}}
    + \lambda_3 L_{\text{pred}} + \lambda_4 L_{\text{aux}}
\]

where:

\[
  L_{\text{outer}} = \text{SmoothL1}(\text{maml\_pred}, \text{target})
\]

\[
  L_{\text{base}} = \text{SmoothL1}(\text{fused}, \text{target})
\]

\[
  L_{\text{pred}} = \text{SmoothL1}(\text{model\_pred}, \text{target})
\]

\[
  L_{\text{aux}} = \text{SmoothL1}(\text{aux\_proj}(\text{llm\_embed\_ca}),
    \text{target})
\]

The weights are:

\[
  \lambda_1 = \text{mc\_blend}
\]

\[
  \lambda_2 = (1 - \text{mc\_blend}) \cdot 0.4
\]

\[
  \lambda_3 = (1 - \text{mc\_blend}) \cdot 0.3
\]

\[
  \lambda_4 = 0.2
\]

where $\text{mc\_blend} \in [0.1, 0.5]$ is derived from Monte Carlo
variance:

\[
  \text{mc\_blend} = \max(0.5 - \sigma^2_{\text{MC}} \cdot 0.5, 0.1)
\]

\textbf{Adaptive Blending}: High MC variance (model uncertainty)
increases the weight on the outer (MAML) loss, which is more stable. Low
variance decreases it. The system adapts loss weights based on model
confidence.

\textbf{SmoothL1 Loss}: Uses $\beta = 0.1$, giving $L_2$ near zero and
$L_1$ far from zero. This prevents saturation unlike Cosine+L1 which
fought itself.

\textbf{Auxiliary Loss}: The aux path carries gradient signal to
cross-attention, embed-weight-net, and text-proj, which sit before the
non-differentiable C boundary.

\subsection{Gradient Balancing}

The embed\_weight\_net branch is often weak. Balancing rescales weak
gradients:

\[
  \text{scale} = \min\left( \frac{\text{norm}_{\text{FT}}}{10 \cdot
    \text{norm}_{\text{embed}}}, 5 \right)
\]

where norms are EMA-smoothed to avoid per-epoch pulses.

\[
  p.\text{grad} \leftarrow p.\text{grad} \cdot \text{scale}
\]

All gradients are then clipped globally:

\[
  \|g\|_2 \leq 1.0
\]

\textbf{Rationale}: Prevents any single module from exploding and ensures
balanced learning across branches.

\subsection{Monte Carlo Sampling}

During forward, the model runs in training mode and takes $T$ stochastic
samples to estimate output variance:

\[
  \text{samples} = [\text{model}(x + \epsilon_1), \ldots,
    \text{model}(x + \epsilon_T)]
\]

where $\epsilon_i \sim \mathcal{N}(0, 0.05^2) \cdot r_{\text{explore}}$
if exploration is above threshold, else zeros (dropout only).

\[
  \sigma^2_{\text{MC}} = \text{Var}(\text{samples})
\]

\textbf{Rationale}: Uncertainty estimates inform loss weighting. Without
dropout, variance collapses and weights become uniform. Gaussian noise at
high exploration keeps estimates meaningful.

\subsection{Frozen Skip Logic}

On a frozen (input, target) window, once the loss falls below
$\texttt{\_frozen\_skip\_floor}$ ($\approx 0.15$), the optimizer step
is skipped:

\[
  \text{skip} = \text{(in\_frozen\_window)} \land
    (L_{\text{loss}} < 0.15)
\]

\textbf{Rationale}: Continuing to step a memorized pair drives gradients
to $\sim 10^{-4}$ and wastes progress. Skipping prevents the sawtooth
artifact in loss logs.

\section{Emotional and Identity Updates}

After each backward pass, side-effects update the C-side backend:

\subsection{Novelty Derivation}

From loss history:

\[
  \text{novelty} = \min\left(
    \frac{|L_{\text{current}} - \bar{L}|}{\bar{L} + 10^{-8}}, 1.0
  \right)
\]

where $\bar{L}$ is the mean of the loss history. This is blended with the
backend's reflection:

\[
  \text{novelty} = \max(\text{novelty}, \text{reflect\_novelty},
    \text{reflect\_drift})
\]

\subsection{Satisfaction and Emotion}

Satisfaction combines low loss and improving loss:

\[
  \text{sat} = \min(\max(1.0 - L, 0) + \max(L_{\text{prev}} - L, 0), 1.0)
\]

Emotion type and strength are picked based on loss delta and novelty:

\[
  \text{type} = \begin{cases}
    \text{SURPRISE} & \text{if } \text{novelty} > 0.7 \\
    \text{LOVE} & \text{if } L_{\text{prev}} - L > 0.03 \\
    \text{HATE} & \text{if } L - L_{\text{prev}} > 0.03 \\
    \text{SURPRISE} & \text{otherwise (weak)}
  \end{cases}
\]

\[
  \text{strength} = \begin{cases}
    \min(\text{novelty} \cdot 0.6, 0.6) & \text{type} = \text{SURPRISE} \\
    \min(\Delta L \cdot 3.0, 0.6) & \text{type} = \text{LOVE/HATE} \\
    0.1 & \text{deadband}
  \end{cases}
\]

\textbf{Deadband}: The $0.03$ threshold prevents small wiggles from
triggering extreme emotions. The $0.1$ deadband emotion keeps arousal
decaying instead of being pinned high.

\section{Inference: LarkosRunner}

\subsection{Initialization and Loading}

LarkosRunner reconstructs only forward-path modules and loads trained
weights via \texttt{load\_checkpoint}:

\[
  \text{Runner} = \text{LarkosRunner}(\text{backend}, \text{ckpt\_path})
\]

The C-side backend is reloaded with saved memory and network states:

\[
  \text{backend.load\_memory()}
\]

\[
  \text{backend.load\_network\_states()}
\]

\textbf{Rationale}: Without the trained backend state, cognitive\_fuse
would see empty neurons and memory, and the readout would not reflect
training dynamics. Both halves (PyTorch weights + C-side state) are needed
together; neither is usable alone.

\subsection{Forward Inference}

The runner's \texttt{step} method mirrors \texttt{TrainingLoop.\_forward}
but without:

\begin{itemize}
  \item MC probes (no uncertainty sampling; frozen dropout)
  \item MAML inner loop (single forward pass only)
  \item Auxiliary loss computation (output path only)
\end{itemize}

All computation runs with \texttt{@torch.no\_grad()}, preventing gradient
accumulation and saving memory.

\subsection{Input and Frozen Window Handling}

Since inference has no training schedule, the frozen-window logic is
disabled:

\[
  \text{frozen\_input} = \text{False}
\]

The fused cognitive vector is always fresh (from the live backend) and
re-normalized:

\[
  \text{fused\_cog} = \text{normalize}(\text{cognitive\_fuse}(\ldots))
\]

This mirrors the behavior of a target-refresh epoch in training.

\subsection{Output}

The runner returns a dictionary with:

\begin{itemize}
  \item \texttt{fused}: The final neuron prediction (NumPy array)
  \item \texttt{model\_pred}: Raw LarkosModel output
  \item \texttt{text\_input}: The driver sentence
  \item \texttt{text\_output}: Decoded text from the fused vector
\end{itemize}

\section{Key Design Decisions}

\subsection{Why Detach at the C Boundary}

The C-side cognitive fusion reads non-differentiable backend objects
(neurons, memory, network history). Python cannot build a computation
graph through C code. Detaching prevents error and forces the gradient
signal through the auxiliary loss instead.

\subsection{Why Frozen Input Windows}

Updating the target every epoch while the transformer input drifts (due to
neuron and memory changes) causes the loss to thrash. Pinning both on the
same refresh schedule prevents chasing a moving target.

\subsection{Why Separate Projections for Bands}

Each band carries different information (LLM embedding, neuron state,
memory). Shared projections would force all streams through the same
learned mapping. Separate projections (and stream embeddings) allow each
stream to learn its own representation and let attention decide balance.

\subsection{Why SmoothL1 Loss}

SmoothL1 with $\beta = 0.1$ avoids saturation of cosine losses and the
fighting behavior of Cosine+L1. It is stable and well-behaved for
vector regression.

\subsection{Why Auxiliary Loss}

The cross-attention, embed-weight-net, and text-proj networks sit before
the C boundary. Without a separate loss on their outputs, they would
receive zero gradient (the path to target is broken). The auxiliary loss
regresses the cross-attention embedding to the target dimension and
provides the needed signal.

\section{Scheduler and Optimization}

Training uses:

\begin{itemize}
  \item Optimizer: Adam with $\text{lr} = \text{BASE\_LR}$,
    $\text{weight\_decay} = 10^{-4}$, $\text{eps} = 10^{-6}$
  \item Scheduler: CosineAnnealingLR with
    $T_{\max} = \text{COSINE\_T\_MAX}$,
    $\eta_{\min} = 10^{-5}$
  \item EMA: Exponential moving average of main model weights (used for
    validation and inference)
\end{itemize}

The scheduler steps after each backward pass, ensuring a smooth cosine
decay over the epoch range.

\section{Conclusion}

The Larkos training loop and runner form a unified system:

\begin{itemize}
  \item \textbf{Training}: Full differentiable forward pass,
    multi-component loss, frozen windows, emotional state management
  \item \textbf{Inference}: Fast forward-only pass, loaded weights and
    backend state, no gradient overhead
\end{itemize}

Both share core modules and math but are optimized for their respective
tasks. The design prioritizes stability (frozen windows, gradient
balancing, adaptive loss weights), interpretability (emotional and
identity updates), and efficiency (EMA, frozen sampling).

\end{document}
